\section{Flexibly Fair VAE} \label{sec:method}

We want to learn fair representations that---beyond being useful for predicting many test-time task labels $y$---can be adapted \textit{simply} and \textit{compositionally} for a variety of sensitive attributes settings $a$ after training.
We call this property \textit{flexible fairness}.
Our approach to this problem involves inducing structure in the latent code that allows for easy manipulation.
Specifically, we isolate information about each sensitive attribute to a specific subspace, while ensuring that the latent space factorizes these subspaces independently. 

\paragraph{Notation}
We employ the following notation:
\begin{itemize}
    \item $x \in \mathcal{X}$: a vector of non-sensitive attributes, for example, the pixel values in an image or row of features in a tabular dataset;
    \item $a \in \{0,1\}^{N_a}$: a vector of binary sensitive attributes;
    \item $z \in \R^{N_z}$: non-sensitive subspace of the latent code;
    \item $b \in \R^{N_b}$: sensitive subspace of the latent code\footnote{
    In our experiments we used $N_b = N_a$ (same number of sensitive attributes as sensitive latent dimensions) to model binary sensitive attributes. 
    But categorical or continuous sensitive attributes can also be accommodated. 
    }.
\end{itemize}
For example, we can express the VAE objective in this notation as
\begin{align}% \label{eq:betavae-new-notation}
    L_{\text{VAE}}(p,q) &= \E_{q(z,b|x,a)} \left[ \log p(x,a|z,b) \right] \nonumber \\
    &\quad - D_{KL} \left[ q(z,b|x,a) || p(z,b)  \right]. \nonumber
\end{align}

In learning a flexibly fair representations $[z,b]=f([x,a])$, we aim to satisfy two general properties: \emph{disentanglement} and \emph{predictiveness}.
We say that $[z,b]$ is \emph{disentangled} if its aggregate posterior factorizes as $q(z,b)=q(z)\prod_j q(b_j)$ and is \emph{predictive} if each $b_i$ has high mutual information with the corresponding $a_i$.
Note that under the disentanglement criteria the dimensions of $z$ are free to co-vary together, but must be independent from all sensitive subspaces ${b_j}$.
We have also specified factorization of the latent space in terms of the aggregate posterior $q(z,b)=\E_{p^{\text{data}}(x)} [q(z,b|x)]$, to match the global independence criteria of group fairness.

\paragraph{Desiderata}
We can formally express our desiderata as follows:
\begin{itemize}
    \item $z \perp b_j \medspace \forall \medspace j$ (disentanglement of the non-sensitive and sensitive latent dimensions);
    \item $b_i \perp b_j \medspace \forall \medspace i\neq j$ (disentanglement of the various different sensitive dimensions);
    \item $\text{MI}(a_j,  b_j)$ is large $\forall \medspace j$ (predictiveness of each sensitive dimension);
\end{itemize}
where $\text{MI}(u,v)=\E_{p(u,v)} \log \frac{p(u,v)}{p(u)p(v)}$ represents the mutual information between random vectors $u$ and $v$.
We note that these desiderata differ in two ways from the standard disentanglement criteria.
The predictiveness requirements are stronger: they allow for the injection of external information into the latent representation, requiring the model to structure its latent code to align with that external information.
However, the disentanglement requirement is less restrictive since it allows for correlations between the dimensions of $z$.
Since those are the non-sensitive dimensions, we are not interested in manipulating those at test time, and so we have no need for constraining them.

If we satisfy these criteria, then it is possible to achieve demographic parity with respect to some $a_i$ by simply removing the dimension $b_i$ from the learned representation i.e. use instead $[z,b]\backslash b_i$.
We can alternatively replace $b_i$ with independent noise.
This adaptation procedure is simple and compositional: if we wish to achieve fairness with respect to a conjunction of binary attributes\footnote{
$\wedge$ and $\vee$ represent logical \emph{and} and \emph{or} operations, respectively.
} $a_i \wedge a_j \wedge a_k$, we can simply use the representation $[z,b]\backslash \{b_i, b_j, b_k\}$.

By comparison, while FactorVAE may disentangle dimensions of the aggregate posterior---$q(z)=\prod_j q(z_j)$---it does not automatically satisfy flexible fairness, since the representations are not predictive, and cannot necessarily be easily modified along the attributes of interest.

\paragraph{Distributions}
We propose a variation to the VAE which encourages our desiderata, building on methods for disentanglement and encouraging predictiveness. 
Firstly, we assume assume a variational posterior that factorizes across $z$ and $b$:
\begin{align}
    q(z,b|x) &= q(z|x)q(b|x).
\end{align}
The parameters of these distributions are implemented as neural network outputs, with the encoder network yielding a tuple of parameters for each input: $(\mu_q(x), \Sigma_q(x), \theta_q(x))=\text{Encoder}(x)$.
We then specify $q(z|x)=\mathcal{N}(z|\mu_q(x), \Sigma_q(x))$ and $q(b|x)=\delta(\theta_q(x))$ (i.e., $b$ is non-stochastic)\footnote{
We experimented with several distributions for modeling $b|x$ stochastically, but modeling this uncertainty did not help optimization or downstream evaluation in our experiments.
}.

Secondly, we model reconstruction of $x$ and prediction of $a$ separately using a factorized decoder:
\begin{align}
    p(x,a|z,b)=p(x|z,b)p(a|b)
\end{align}
where $p(x|z,b)$ is the decoder distribution suitably chosen for the inputs $x$, and $p(a|b)=\prod_j\text{Bernoulli}(a_j|\sigma(b_j))$ is a factorized binary classifier that uses $b_j$ as the logit for predicting $a_j$ ($\sigma(\cdot)$ represents the sigmoid function).
Note that the $p(a|b)$ factor of the decoder requires no extra parameters.

Finally, we specify a factorized prior $p(z,b)=p(z)p(b)$ with $p(z)$ as a standard Gaussian and $p(b)$ as Uniform.

\paragraph{Learning Objective}
Using the encoder and decoder as defined above, we present our final objective:
\begin{align} \label{eq:ffvae}
    \nonumber L_{\text{FFVAE}}(p,q) &= \E_{q(z,b|x)} [ \log p(x|z,b) + \alpha \log p(a|b)] \nonumber \\
    &\quad\quad - \gamma D_{KL}(q(z,b)||q(z)\prod_j q(b_j))
     \nonumber\\
    &\quad\quad - D_{KL} \left[ q(z,b|x) || p(z,b) \right].
\end{align}
It comprises the following four terms, respectively:
a \textit{reconstruction} term which rewards the model for successfully modeling non-sensitive observations;
a \textit{predictiveness} term which rewards the model for aligning the correct latent components with the sensitive attributes; 
a \textit{disentanglement} term which rewards the model for decorrelating the latent dimensions of $b$ from each other and $z$;
and a \textit{dimension-wise KL} term which rewards the model for matching the prior in the latent variables.
We call our model \mbox{FFVAE} for Flexibly Fair VAE (see Figure \ref{modelpic} for a schematic representation).

The hyperparameters $\alpha$ and $\gamma$ control aspects relevant to flexible fairness of the representation. $\alpha$ controls the alignment of each $a_j$ to its corresponding $b_j$ (predictiveness), whereas $\gamma$ controls the aggregate independence in the latent code (disentanglement).

The $\gamma$-weighted total correlation term is realized by training a binary adversary  to approximate the log density ratio $\log \frac{q(z,b)}{q(z)\prod_j q(b_j)}$.
The adversary attempts to classify between ``true'' samples from the aggregate posterior $q(z,b)$ and ``fake'' samples from the product of the marginals $q(z)\prod_j q(b_j)$ (see Appendix \ref{sec:disc_approx} for further details).
If a strong adversary can do no better than random chance, then the desired independence property has been achieved.

We note that our model requires the sensitive attributes $a$ at training time but not at test time.
This is advantageous, since often these attributes can be difficult to collect from users, due to practical and legal restrictions, particularly for sensitive information \citep{elliot2008new,division1983}.
